% TODO make hw stylesheet

\documentclass[10pt]{article}
\usepackage[letterpaper,top=1in,left=1in,right=1in]{geometry}
\usepackage[dvipsnames,svgnames]{xcolor}
\usepackage[shortlabels]{enumitem}
\usepackage[framemethod=TikZ]{mdframed}
\usepackage{amsmath,amssymb,amsthm}
\usepackage[colorlinks]{hyperref}
\usepackage{multicol}
\usepackage{tabularx}
\usepackage{bbm}
\usepackage{tikz-cd}

\hypersetup{
    colorlinks=true,
    linkcolor=blue,
    filecolor=magenta,
    urlcolor=magenta,
    pdftitle={MAT 535 HW}
}

\usepackage{mathtools}
\usepackage{thmtools}
\usepackage[textsize=footnotesize]{todonotes}
\usepackage{titling}
\usepackage{cancel}

\reversemarginpar
\mdfdefinestyle{mdbluebox}{roundcorner=10pt,innerbottommargin=9pt,
linecolor=blue,backgroundcolor=TealBlue!5,}
\declaretheoremstyle[headfont=\sffamily\bfseries\color{MidnightBlue},
mdframed={style=mdbluebox},]{thmbluebox}

\mdfdefinestyle{mdbloobox}{frametitlefont=\bfseries,innerbottommargin=8pt,
nobreak=true,backgroundcolor=TealBlue!5,linecolor=blue,}
\declaretheoremstyle[headfont=\bfseries\color{MidnightBlue},
mdframed={style=mdbloobox},headpunct={\\[3pt]},postheadspace=0pt,]{thmbloobox}

\mdfdefinestyle{mdredbox}{frametitlefont=\bfseries,innerbottommargin=8pt,
nobreak=true,backgroundcolor=Salmon!5,linecolor=RawSienna,}
\declaretheoremstyle[headfont=\bfseries\color{RawSienna},
mdframed={style=mdredbox},headpunct={\\[3pt]},postheadspace=0pt,]{thmredbox}

\mdfdefinestyle{mdgreenbox}{linecolor=ForestGreen,backgroundcolor=ForestGreen!5,
linewidth=2pt,rightline=false,leftline=true,topline=false,bottomline=false,}
\declaretheoremstyle[headfont=\bfseries\sffamily\color{ForestGreen!70!black},
mdframed={style=mdgreenbox},headpunct={ --- },]{thmgreenbox}

\mdfdefinestyle{mdorangebox}{linecolor=RedOrange,backgroundcolor=RedOrange!5,
linewidth=2pt,rightline=false,leftline=true,topline=false,bottomline=false,}
\declaretheoremstyle[headfont=\bfseries\sffamily\color{RedOrange!70!black},
mdframed={style=mdorangebox},headpunct={ --- },]{thmorangebox}

\mdfdefinestyle{mdblackbox}{linecolor=black,backgroundcolor=RedViolet!5!gray!5,
linewidth=3pt,rightline=false,leftline=true,topline=false,bottomline=false,}
\declaretheoremstyle[mdframed={style=mdblackbox}]{thmblackbox}

\declaretheoremstyle[spaceabove=6pt,spacebelow=6pt]{boldhead}
\declaretheoremstyle[spaceabove=6pt,spacebelow=6pt,headfont=\normalfont\itshape]{italichead}

\declaretheorem[style=boldhead,name=Problem]{problem}
\declaretheorem[style=boldhead,name=Problem,numbered=no]{problem*}
\declaretheorem[style=boldhead,name=Setup]{setup} % for config geo
\declaretheorem[style=boldhead,name=Example,sibling=problem]{example}
\declaretheorem[style=boldhead,name=$\bigstar$ Problem,sibling=problem]{niceprob}
\declaretheorem[style=boldhead,name=Theorem,sibling=problem]{theorem}
\declaretheorem[style=boldhead,name=Corollary,sibling=theorem]{corollary}
\declaretheorem[style=boldhead,name=Proposition,sibling=theorem]{prop}
\declaretheorem[style=boldhead,name=Lemma,numberwithin=problem]{lemma}
\declaretheorem[style=boldhead,name=Theorem,numbered=no]{theorem*}
\declaretheorem[style=boldhead,name=Lemma,numbered=no]{lemma*}
\declaretheorem[style=boldhead,name=Claim,numberwithin=problem]{claim}
\declaretheorem[style=boldhead,name=Claim,numbered=no]{claim*}
\declaretheorem[style=boldhead,name=Remark,numberwithin=problem]{remark}
\declaretheorem[style=boldhead,name=Remark,numbered=no]{remark*}
\declaretheorem[style=boldhead,name=Definition,sibling=theorem]{definition}
\declaretheorem[style=boldhead,name=Definition,numbered=no]{definition*}

\providecommand{\ol}{\overline}
\providecommand{\tensor}{\otimes}
\renewcommand{\hom}{\operatorname{Hom}}
\providecommand{\tor}{\operatorname{Tor}}
\providecommand{\Gal}{\operatorname{Gal}}
\providecommand{\ceil}[1]{\left\lceil #1 \right\rceil}
\providecommand{\floor}[1]{\left\lfloor #1 \right\rfloor}
\newcommand{\dd}{\,\mathrm{d}}
\providecommand{\ul}{\underline}
\providecommand{\eps}{\varepsilon}
\providecommand{\half}{\frac{1}{2}}
\providecommand{\CC}{\mathbb C}
\providecommand{\FF}{\mathbb F}
\providecommand{\II}{\mathbb I}
\providecommand{\NN}{\mathbb N}
\providecommand{\QQ}{\mathbb Q}
\providecommand{\RR}{\mathbb R}
\providecommand{\ZZ}{\mathbb Z}
\providecommand{\ind}{\mathbbm 1}
\providecommand{\dg}{^\circ}
\providecommand{\ii}{\item}
\providecommand{\setdiff}{\smallsetminus}
\providecommand{\alert}[1]{{\sffamily\textbf{\textcolor{blue}{#1}}}}
\providecommand{\inv}{^{-1}}
\providecommand{\mb}[1]{\mathbf{#1}}
\newcommand{\what}{\widehat}

\newenvironment{soln}{\begin{proof}[Solution]}{\end{proof}}

\providecommand{\pow}{\operatorname{Pow}}
\providecommand{\vocab}[1]{{\textbf{\textcolor{ForestGreen}{#1}}}}
\providecommand{\lcm}{\operatorname{lcm}}
\providecommand{\abs}[1]{\left\lvert #1\right\rvert}
\providecommand{\smabs}[1]{\lvert #1\rvert}
\providecommand{\innerprod}[1]{\langle\!\langle #1\rangle\!\rangle}
\providecommand{\norm}[1]{\left\| #1\right\|}
\providecommand{\smnorm}[1]{\| #1\|}
\providecommand{\gr}{\operatorname{gr}}

\usepackage{mathrsfs}

% place title at top right
\setlength{\droptitle}{-8ex}
\pretitle{\begin{flushright}\Large\bfseries}
\posttitle{\par\end{flushright}}
\preauthor{\begin{flushright}}
\postauthor{\end{flushright}}
\predate{\begin{flushright}}
\postdate{\end{flushright}}

\title{MAT 535 HW07}
\author{\large Aditya Pahuja}
\date{\large Due 04/02}
\begin{document}
\maketitle

\begin{problem}
    [\S13.6\#5]
    Prove that there are finitely many roots of unity
    in any finite extension $K$ of $\QQ$.
\end{problem}

\begin{soln}
    Let $L$ be the Galois closure of $K$ over $\QQ$,
    which is also a finite extension of $\QQ$.

    \begin{lemma}
        For each positive integer $k$, $\varphi(n) = k$
	has finitely many positive integer solutions $n$.
    \end{lemma}
    \begin{proof}
	Let $p$ be a prime number and $e$ a positive integer.
	If $p > 3$, then
	\begin{equation*}
	    \varphi(p^e) = p^{e-1}(p-1) > \frac{p^e}2 > \sqrt{p^e},
	\end{equation*}
	and if $p \le 3$ then $\varphi(p^e) \ge\frac{\sqrt{p^e}}2$ instead.
	Thus, for an arbitrary positive integer $n$
	with prime factorization $n = p_1^{e_1}\cdots p_r^{e_r}$
	for distinct primes $p_j$ and positive integers $e_j$,
	\begin{equation*}
	    \varphi(n) = \prod_{j=1}^r \varphi(p_j^{e_j})
	    \ge \frac 14\prod_{j=1}^n \sqrt{p_j^{e_j}} = \frac{\sqrt n}4.
	\end{equation*}
	In other words, if $\varphi(n) = k$, then
	$k\ge\frac{\sqrt n}4$ implies $n \le 16k^2$,
	and there are only finitely many such $n$.
	(The constant $16$ can be improved to $2$
	with a little bit more care
	when analyzing $p = 2,3$;
	in particular $\varphi(3^e) > \sqrt{3^e}$ for all $e$
	and $\varphi(2^e) \ge \sqrt{2^e}$ when $e > 1$,
	so the only obstruction to a lower bound of $\varphi(n)\ge \sqrt n$     
	is $\varphi(2) = 1 < \sqrt 2$,
	giving the inequality $\varphi(n) \ge \sqrt{\frac n2}$
	which is sharp if and only if $n = 2$.)
    \end{proof}

    Each root of unity $\zeta_m\in K$ with multiplicative order $m$
    generates the Galois subextension $\QQ(\zeta_m)$ of $L$.
    By the fundamental theorem of Galois theory, $\Gal(\QQ(\zeta_m)/\QQ)$,
    which has order $\varphi(m)$,
    is a quotient of the finite group $\Gal(L/\QQ)$,
    so $\varphi(m)$ must divide its order $N$.
    By the lemma, only finitely many positive integers $m$
    have the property that $\varphi(m)$ divides $N$,
    so only finitely many roots of unity can be in $K$.
    (This also yields a crude upper bound of
    \begin{equation*}
	\sum_{\substack{m\in\NN\\ \varphi(m)\mid N}} \varphi(m)
    \end{equation*}
    for the number of roots of unity in $K$,
    since $\zeta_m$ has $\varphi(m)$ Galois conjugates.)
\end{soln}

\begin{problem}
    [\S13.6\#7]
    Use the M\"obius inversion formula to show that
    \begin{equation*}
	\Phi_m(x) = \prod_{d\mid n}(x^d - 1)^{\mu(n/d)}.
    \end{equation*}
\end{problem}

\begin{proof}
    The M\"obius inversion formula in its usual (additive) form says that
    for functions $f,F\colon\NN\to\CC$,
    \begin{equation*}
	F(n) = \sum_{d\mid n}f(d) \iff
	f(n) = \sum_{d\mid n}F(d)\mu(n/d).
    \end{equation*}
    By exponentiating, we get the identity 
    \begin{equation*}
	F(n) = \prod_{d\mid n} f(d)\iff
	f(n) = \prod_{d\mid n} F(d)^{\mu(n/d)},
    \end{equation*}
    which gives the desired equation
    by taking $F(n) = x^n - 1$ and $f(n) = \Phi_n(x)$ since, by definition,
    \begin{equation*}
	x^n - 1 = \prod_{d\mid n}\Phi_d(x).\qedhere
    \end{equation*}
\end{proof}

\pagebreak

\begin{problem}
    [\S14.3\#2]
    Write out the multiplication tables for $\FF_4$ and $\FF_8$.
\end{problem}

\begin{soln}
    The group of units of $\FF_4$ is cyclic of order $3$,
    so write $\FF_4 = \{0, 1, a, a^2\}$ where $a$
    generates the group of units. Then the multiplication table is
    \begin{center}
	\begin{tabular}{c | cccc}
	    & $0$ & $1$ & $a$ & $a^2$\\
	    \hline
	    $0$ & $0$ & $0$ & $0$ & $0$\\
	    $1$ & $0$ & $1$ & $a$ & $a^2$\\
	    $a$ & $0$ & $a$ & $a^2$ & $1$\\
	    $a^2$ & $0$ & $a^2$ & $1$ & $a$
	\end{tabular}
    \end{center}
    Similarly, the group of units of $\FF_8$ is cyclic of order $7$,
    so write $\FF_8 = \{0, 1, b, \dots, b^6\}$
    where $b$ generates the group of units.
    Then the multiplication table is
    \begin{center}
	\begin{tabular}{c|cccccccc}
	    &$0$&$1$&$b$&$b^2$&$b^3$&$b^4$&$b^5$&$b^6$\\
	    \hline
	    $0$&$0$&$0$&$0$&$0$&$0$&$0$&$0$&$0$\\
	    $1$&$0$&$1$&$b$&$b^2$&$b^3$&$b^4$&$b^5$&$b^6$\\
	    $b$&$0$&$b$&$b^2$&$b^3$&$b^4$&$b^5$&$b^6$&$1$\\
	    $b^2$&$0$&$b^2$&$b^3$&$b^4$&$b^5$&$b^6$&$1$&$b$\\
	    $b^3$&$0$&$b^3$&$b^4$&$b^5$&$b^6$&$1$&$b$&$b^2$\\
	    $b^4$&$0$&$b^4$&$b^5$&$b^6$&$1$&$b$&$b^2$&$b^3$\\
	    $b^5$&$0$&$b^5$&$b^6$&$1$&$b$&$b^2$&$b^3$&$b^4$\\
	    $b^6$&$0$&$b^6$&$1$&$b$&$b^2$&$b^3$&$b^4$&$b^5$\\
        \end{tabular}
    \end{center}
\end{soln}

\begin{problem}
    [\S14.3\#8]
    Determine the splitting field for the polynomial $x^p - x - a$
    over $\FF_p$ where $a\in\FF_p^\times$. Show explicitly
    that the Galois group is cyclic.
\end{problem}

\begin{proof}
    We have shown on previous homework that if $f(x) = x^p - x - a$
    does not have any roots in a given field $K$ of characteristic $p$
    then it is irreducible, and that if $f(x)$ has any root
    $\alpha\in K$ then it splits in $K$,
    with the roots of $f(x)$ being $\alpha + j$, $j\in\FF_p$.

    Since $y^p - y = 0\ne a$ for any $y\in\FF_p$,
    $f(x)$ is irreducible over $\FF_p[x]$.
    Then $K = \FF_p[x]/(f(x))$ contains a root $\ol x$
    (the image of $x$ under the projection)
    of $f(x)$, so $f(x)$ splits in this field.
    Moreover $K = \FF_p(\ol x)$ is generated by
    the roots of $x^p - x - a$, so $K$ is the splitting field of $f(x)$.
    By construction, $[K : F] = \deg(f) = p$, hence $K = \FF_{p^p}$.
    We noted above that the roots of $f(x)$ are
    $\ol x + j$ for $j\in\FF_p$.
    I claim that the maps $\varphi_j\colon K\to K$ defined by
    \begin{equation*}
	\varphi_j(a_0 + a_1\ol x + \cdots + a_{p-1}\ol x^{p-1})
	= a_0 + a_1(\ol x + j) + \cdots + a_{p-1}(\ol x + j)^{p-1}
    \end{equation*}
    for $a_0,\dots,a_{p-1}\in\FF_p$ are the automorphisms of $K$.
    To see this, observe that
    \begin{align*}
	\varphi_a(a_0 + a_1\ol x + \cdots + a_{p-1}\ol x^{p-1})
	&= a_0 + a_1(\ol x + a) + \cdots + a_{p-1}(\ol x + a)^{p-1}\\
	&= a_0 + a_1\ol x^p + \cdots + a_{p-1}\ol x^{p(p-1)}\\
	&= (a_0 + a_1\ol x + \cdots + a_{p-1}\ol x^{p-1})^p,
    \end{align*}
    which is the Frobenius automorphism,
    and for any other $j\in\FF_p$,
    we can write $j = na$ for some $n\in\FF_p$
    to get $\varphi_j = \varphi_a^n$, so that
    $\varphi_a$ generates a cyclic subgroup of $\Gal(K/\FF_p)$
    of order $p$. Since $[K : \FF_p] = p$, this implies
    $\Gal(K/\FF_p)$ is equal to $\langle \varphi_a\rangle$.
\end{proof}










\end{document}
